%% EAAI first draft converted to Elsevier elsarticle format.
%% Target journal: Engineering Applications of Artificial Intelligence.
\documentclass[preprint,review,12pt]{elsarticle}

\usepackage[utf8]{inputenc}
\usepackage[T1]{fontenc}
\usepackage{amsmath,amssymb,amsfonts}
\usepackage{graphicx}
\usepackage{booktabs}
\usepackage{tabularx}
\usepackage{array}
\usepackage{algorithm}
\usepackage{algpseudocode}
\usepackage{lineno}
\usepackage{xcolor}
\usepackage{newunicodechar}
\newunicodechar{▹}{$\triangleright$}
\newunicodechar{ï}{\"i}
\newunicodechar{Γ}{\Gamma}
\newunicodechar{’}{'}
\newunicodechar{“}{``}
\newunicodechar{”}{''}
\newunicodechar{–}{--}
\newunicodechar{—}{---}
\newunicodechar{−}{-}
\newunicodechar{×}{\times}
\graphicspath{{figures/}}

\bibliographystyle{elsarticle-num}

\journal{Engineering Applications of Artificial Intelligence}

\begin{document}
\begin{frontmatter}

\title{Graph-Aware Value Learning for Support-Hour Planning in Synthetic Personalized Medicine Manufacturing Networks}

\author[inst1]{Zhaowei Li\corref{cor1}}
\ead{zli686@gatech.edu}
\cortext[cor1]{Corresponding author.}
\author[inst1]{Chin-Yuan Tseng}
% \ead{...}  % add publication email if you want it listed for the second author
\affiliation[inst1]{organization={H. Milton Stewart School of Industrial and Systems Engineering, Georgia Institute of Technology},
            city={Atlanta},
            state={GA},
            country={USA}}

\begin{abstract}
Advance booking of support labor couples current expenditure with future
manufacturing readiness and modeled patient eligibility. We study this
decision in a four-site synthetic, patient-indexed manufacturing network.
A graph-convolutional temporal-difference value model augments an eight-step
model predictive controller (MPC), while the base heuristic retains routing,
purchasing, and resource-sharing decisions. We distinguish improvement over
fixed support allocation from improvement over strong planning. In a
five-family study, neither development-selected finalist passed the
prespecified H8-MPC screen: Value-TD saved 0.85\% mean cost with an interval
crossing zero, and soft actor--critic cost 3.50\% more. A post hoc comparison
on the same 120 evaluated worlds found 2.16\% lower mean cost and 5.675
fewer simulated patient losses per world for Value-TD than for MDL-2 with
fixed additional support. A subsequent 72-world development comparison
used that fixed-support policy as its prospective primary baseline and H8
as a strong reference. Original Value-TD saved 2.34\% and reduced mean
losses by 5.181 per world. A fixed half-residual variant failed its
condition-wise patient guardrail and was not adopted. These findings
support a bounded cost--service comparison, not independent confirmation,
isolated graph or learning benefits, or clinical safety. External
calibration and fresh confirmation remain necessary.
\end{abstract}

\begin{keyword}
Graph convolutional networks \sep Temporal-difference learning \sep
Model predictive control \sep Capacity planning \sep Synthetic simulation
\sep Personalized manufacturing
\end{keyword}
\end{frontmatter}
\linenumbers

\begin{center}
\textbf{Working draft:} evidence and source checks by accountable human
authors, final declarations, and submission approval remain pending.
\end{center}

\section{Introduction}
\label{sec:introduction}
Personalized manufacturing links individual demand to non-interchangeable
products. In the simulation considered here, a specimen and its resulting
product retain one patient's identity, while eligibility, specimen viability,
and inventories evolve over time. Qualified support labor prepares work for
manufacturing but does not shorten biological production or release testing.
Published process-cost and labor studies motivate separating these resource
requirements; they do not calibrate the staffing response of this
simulator~\cite{Ran2020DecentralizedCART,Fitzgerald2022AutomatedASC}.

The operational question is whether state-dependent advance booking improves
on a fixed support schedule while protecting simulated service outcomes.
A second, distinct question is whether learning adds value beyond a competent
planner using the same public observations. A fixed schedule and a planner
are not interchangeable baselines: they use different information and
computation, and may trade expenditure against patient losses differently.

We investigate Value-TD, a graph-aware terminal-cost estimator coupled to
eight-step MPC. The planner chooses support hours; the neural model does
not directly output actions. Graph convolution, residual learning, and
learned terminal values are established ideas
\cite{Kipf2017GCN,Johannink2019ResidualRL,Hansen2022TDMPC}.
The contribution considered here is their task-specific integration and
evaluation under delayed labor commitments, partial public observations,
and patient-indexed service consequences. We do not claim a new
reinforcement-learning backbone, a first-in-literature combination, or a
general superiority result.

The paper makes three bounded contributions.
\begin{itemize}
\item It specifies a synthetic support-hour decision problem with common
base operations, delayed paid commitments, full settlement accounting,
and an explicit public-information interface.
\item It describes an observed-trajectory TD value correction to a fixed
short-horizon planner, separating the learned score from the operational
heuristic, physical dynamics, and actor--critic alternatives.
\item It reports a staged comparison that retains a failed strong-planner
screen, a post hoc fixed-support supplement, and a prospectively specified
development test that rejected residual attenuation. Cost, service,
computation, and unresolved component attribution are reported together.
\end{itemize}

Future primary comparisons use MDL-2 with fixed additional support, with H8-MPC
retained as a strong heuristic reference. This prospective choice does not
replace the original five-family study's H8 endpoint. Favorable fixed-support
comparisons do not rescue that failed screen. Earlier full-routing studies
address a different action space and are retained in the separate
\texttt{historical-research-record.tex}, not pooled into this paper's
support-hour leaderboard. All reported losses and costs are synthetic;
publication readiness and operational validity are not established.

\section{Support-Hour Decision Problem}
\label{sec:problem}
\subsection{Operational boundary}
The task contains four manufacturing sites on a fixed network. Each world
has 48 control epochs followed by 16 settlement epochs. During control,
the policy books additional support hours by site. Bookings arrive after
two epochs; there are no new bookings during settlement, but outstanding
commitments and their consequences continue to be accounted for. Ordinary
staffing is four hours per site and epoch. Additional support acts on the
simulated preparatory-work requirement; it does not directly speed cell
growth, remove equipment constraints, or change patient deterioration.
The configuration labels epochs as modeled weeks, but that label is not
validated against a process calendar. A single support pool is assumed
qualified at all sites; individual credentials, shifts, and staff travel
are not modeled. Flexible requests are capped at four hours per site and
eight hours network-wide, with unused budget allowed.

All roles use the same MDL-2 base operations for routing, purchasing, and
resource sharing. In this task, MDL-2 denotes the two-period mean-demand
look-ahead heuristic, not a minimum-description-length algorithm. The
learned intervention changes support allocation only. It does not inherit
the action space or empirical status of historical routing/DDPG studies.
During settlement, the executed purchasing rule uses all waiting patient
identities, net of available reagents and outstanding orders, subject to
the common cutoff. Inherited proposal text instead described support-ready
patients only. We report the executed all-waiting rule; this documentation
discrepancy was not silently repaired. The same implementation applies to
all roles, although their different states produce different purchases.

The public view contains currently available patient and operational records,
commitments, and observation-derived beliefs. Exact latent work, response
states, future events, and future random tapes are not controller inputs.
The public predictor and response filter approximate these quantities.
Consequently the feature representation is a compressed information state,
not a proven sufficient statistic for a fully observed Markov decision process.

\subsection{Objective and observations}
Let $c_t$ be the full recorded synthetic cost, including operational,
support, patient-loss, expiry, urgency, and terminal-accounting terms.
The finite-horizon objective is
\begin{equation}
J(\pi)=\mathbb{E}_{\pi}\!\left[\sum_{t=0}^{63}c_t\right].
\label{eq:support-objective}
\end{equation}
Discounting is one. Costs remain synthetic weighted objective units, not
site-calibrated monetary expenditure. Simulated patient loss is a modeled
service outcome and must not be interpreted as a measured death or a
clinical treatment effect.
In particular, the recorded loss event includes modeled non-delivery from
ineligibility or material/product loss, and the field named expiry cost
also includes in-process discard. These accounting names do not identify
clinical mortality or independent biological batch-failure probabilities.

The shared value representation has 31 features per site: patient counts
and urgency, age and eligibility summaries, inferred remaining-work bounds,
resources and transit pipelines, support commitments and previous booking,
estimated response, production stages, time remaining, and demand.
Observed and forecast states use the same schema. A shared schema does
not eliminate forecast error or forecast-to-observed distribution shift.

The unchanged scenario set comprises no capacity change, persistent
unannounced change, and fast independent fluctuation. All three receive
equal weight in the reported aggregates. These are author-selected
synthetic conditions, not estimated frequencies of real manufacturing regimes.

\subsection{Comparator roles and engineering assumptions}
The prospective primary policy, MDL-2 fixed support, books
$[2,2,2,2]$ additional hours at every control epoch. Thus each full world
books and receives 384 additional hours. It retains normal staffing,
delivery lag, feasibility, and charges. It requires no neural value,
response filter, or MPC rollout. H8-MPC instead allocates support through
eight-step public-model planning; H16-MPC is a separately charged
longer-horizon reference where included. Actual support usage is not
equalized across these whole-controller comparisons.

\begin{table}[tbp]
\centering
\caption{Engineering interpretation of the support-hour task. None of the
listed numerical choices is established as an externally calibrated site input.}
\label{tab:engineering-scope}
\small
\begin{tabularx}{\textwidth}{p{0.25\textwidth}XX}
\toprule
Element & Implemented interpretation & Evidence boundary \\
\midrule
Time and commitments & 48 control plus 16 settlement epochs; two-epoch
support lag & Modeled-week label, not a validated process calendar. \\
Staffing & Four ordinary hours/site/epoch; fixed comparator adds two &
Synthetic allocation rule, not an optimized or measured staffing standard. \\
Work response & Support advances preparatory readiness under equipment
and availability constraints & No empirical dose--response curve for extra labor. \\
Patient losses & Simulated eligibility and lifecycle outcomes & Not
observed clinical mortality or infusion benefit. \\
Cost & Weighted operational and patient-facing objective & Not current
currency-valued realized savings. \\
\bottomrule
\end{tabularx}
\end{table}

External process studies contextualize costs, qualification, and task duration
\cite{Ran2020DecentralizedCART,Fitzgerald2022AutomatedASC}. Their products,
sites, units, and accounting boundaries differ from this model. Published
batch-failure assumptions and patient-pathway summaries
\cite{Lopes2020VeinToVein,Nastoupil2020StandardOfCare} do not identify
this simulator's response law or deterioration process. The external
operational-calibration gap therefore remains open.

\section{Related Methods and Claim Boundary}
\label{sec:literature}
Graph convolution supplies neighborhood aggregation, not a new learning
principle~\cite{Kipf2017GCN}. Residual reinforcement learning can combine
a conventional controller's action with a learned action
correction~\cite{Johannink2019ResidualRL}. Our current residual instead
modifies a terminal cost used to rank candidate support allocations.
This distinction separates the present Value-TD controller from the
historical residual DDPG action controller.

TD-MPC jointly learns task-oriented latent dynamics and a terminal value
for short-horizon planning~\cite{Hansen2022TDMPC}. Here the public predictor
is fixed, the value model receives graph-structured operational features,
and learning uses complete observed simulation trajectories. We neither
implement full TD-MPC nor claim to introduce learned terminal costs.
The limited positioning claim is an application-specific formulation and
implementation, not an exhaustive novelty or priority claim.

Comparisons must also separate complete-controller performance from component
attribution. A graph/self-only contrast tests neural message adjacency while
preserving the physical network. A learned/initial contrast tests the value
of training relative to the declared initializer. Cost gains over fixed
staffing alone establish neither contribution. This separation follows the
need for seed-aware, uncertainty-aware evaluation in deep reinforcement
learning~\cite{Henderson2018DRLMatters,Agarwal2021RLiable}.

\section{Observed-Trajectory Value-TD Planning}
\label{sec:method}
\subsection{Planner and graph value}
Let $b_t$ denote the public information state and $H(b)$ the fixed terminal
patient-risk heuristic. The value model is
\begin{equation}
V_\theta(b)=H(b)+s f_\theta(b), \qquad s=10^6.
\label{eq:residual-value}
\end{equation}
Two graph-convolution layers of width 32 encode the four sites, followed by
mean pooling and a scalar residual head. The final linear head is initialized
to zero. The graph and self-only architectures use the same trainable tensors
at initialization; self-only replaces message adjacency by the identity
matrix without altering operational links.

For each of 16 candidate action sequences, the planner simulates eight
public-model steps under three response quantiles. With fixed weights $w_q$,
it minimizes
\begin{equation}
S_\theta(a;b_t)=
\sum_{q=1}^{3}w_q\left[
\sum_{k=0}^{7}\widehat c_{t+k}(a,q)
+H(\widehat b_{t+8}^{a,q})
+s f_\theta(\widehat b_{t+8}^{a,q})
\right].
\label{eq:current-value-mpc-score}
\end{equation}
It executes the first support action and replans after the next public
observation. Both terminal components are zero at the fully settled boundary.
Plain H8 uses the same candidates, predictor, and heuristic with no learned
residual. A zero-residual initializer therefore recovers H8 scores;
this is an implementation identity, not a newly estimated performance claim.
H16 changes the planning horizon and is not an equal-computation treatment.

\subsection{Training as performed}
The five-family study's Value-TD arm started from zero residuals rather than
the historical native-tail model. For each complete 64-transition trajectory,
let $u=\min(t+8,64)$. The residual target is
\begin{equation}
y_t=
\frac{\sum_{k=t}^{u-1}c_k-H(b_t)}{s}
+\mathbf{1}_{u<64}\left[
\frac{H(b_u)}{s}+f_{\theta^-}(b_u)\right],
\label{eq:current-observed-td}
\end{equation}
where $\theta^-$ is the model before that trajectory's fitting boundary.
Targets are fixed for 32 minibatch updates of size 64; samples are drawn
with replacement from the trajectory's 64 starting states. The learner
uses Adam and mean-squared target error. Each training block contains
96 worlds. The selected Value-TD learning rate is $10^{-4}$; the family
selection procedure and fresh-block comparisons are described below.

The full settlement costs enter these observed returns, with zero terminal
bootstrap. There are no native branches or forecast-tail training labels
in this recipe. Model-based candidate forecasts remain part of action
selection and are charged separately. This is value-based learning coupled
to planning, not TD3, DDPG actor optimization, teacher imitation, or
deployment-time parameter adaptation. Evaluation weights remain frozen.

A conceptual limitation remains: the learned value describes observed
states under the current behavior, but is queried at forecast candidate
end states. Approximation error and distribution shift can change rankings.
Neither lower TD loss nor actual optimizer execution proves improved
operational decisions; the endpoint controls must supply that evidence.

\subsection{Evidence stages and estimands}
\begin{table}[tbp]
\centering
\caption{Distinct support-hour evidence stages. Sample counts refer to paired
evaluation worlds, not independent real patients. Historical endpoints remain unchanged.}
\label{tab:evidence-stages}
\small
\begin{tabularx}{\textwidth}{p{0.24\textwidth}p{0.24\textwidth}X}
\toprule
Stage & Comparator and worlds & Permitted inference \\
\midrule
Five-family selection and fresh training & H8 primary; 120 new test worlds
across five fresh training blocks & Frozen family screen and matched graph/learning
contrasts; no qualifying winner. \\
Fixed-support supplement & Added MDL-2 fixed support on those same 120 worlds &
Post hoc whole-controller cost--service comparison, not new confirmation. \\
Fixed attenuation development & MDL-2 fixed support primary, H8 reference;
72 fresh worlds with three old development models & Prospective bounded
variant decision conditional on those models; not fresh-training confirmation. \\
\bottomrule
\end{tabularx}
\end{table}

For costs, positive saving means reference minus candidate; for losses,
positive excess means candidate minus reference. Percentage savings divide
the difference of overall means by the reference mean. In each reported
stage, 2,000 descriptive bootstrap draws resample model/training blocks
commonly across conditions and worlds independently within each condition.
A common replicate index across different conditions is not a matched patient.
Intervals are descriptive, not multiplicity-adjusted formal tests or
guarantees of safety. Raw world, block, and condition harms remain available.

\section{Computational Experiments}
\label{sec:experiments}
The following sections preserve the order of evidence generation. Switching
the future primary baseline to fixed support does not retrospectively
change the five-family protocol or erase its failed H8 screen.

\subsection{Five-family fixed-benchmark comparison: neither finalist passes}
\label{sec:family-selection-development}
The prospective family study compared DDPG, TD3, SAC, PPO, and observed-TD8
value-MPC from scratch on the same synthetic support-hour task. All families
used the same 31 public features, unchanged reward and physical dynamics,
three capacity conditions, and MDL-2 base operations. Plain H8-MPC remained
the primary operational reference; H16-MPC was the secondary compute reference.
This restricted task is not pooled with the routing-primary results.
In particular, the value-TD arm learned a zero-initialized residual from its
own complete observed trajectories; it did not reproduce the earlier
native-branch or forecast-tail recipes.

Each family used learning rates $10^{-4}$ and $3\times10^{-4}$ in three
development blocks, with 96 training worlds per configuration and block.
The 2,880 training trajectories were followed by 864 frozen development
evaluations on 24 shared worlds per block under 12 roles. A frozen
patient-first, then cost selection rule chose one learning rate per family
and two distinct families: value-TD at $10^{-4}$ and SAC at
$3\times10^{-4}$. These were development selections, not winner declarations.
Their development mean cost savings against H8 were 1.63\% and $-1.39\%$,
with 12.24 and 5.38 fewer modeled patient losses per world, respectively.

Five new blocks then trained both finalists with graph and self-only
adjacency, yielding another 1,920 trajectories. Matched graph/self-only
models shared initial trainable tensors and training seeds; only adjacency
differed, while subsequent endogenous trajectories could differ. All 20
final models and the graph initial controls were sealed before independent
evaluation. Each new block supplied 24 test worlds for nine frozen roles:
the two finalists' graph-final, self-only-final, and graph-initial models,
H8, H16, and the predeclared historical conditional graph-TD reference.
The historical models were evaluation-only. All 1,080 independent evaluations
completed with zero optimizer updates. The package contained 6,744 main
trajectories and no clones or old-holdout reuse.

Every world included 48 decisions and 16 settlement transitions. Training
used 32 minibatch iterations of size 64 after each world. Actor--critic
targets used $\gamma=1$, reward equal to negative full cost divided by
$10^6$, and all settlement costs in the final control transition. Off-policy
replay was capped at 8,192 transitions; PPO used the current episode and
latent-action log probabilities, while SAC entropy concerned the proposal
before projection. Extra critics, planning, and sample reuse were charged,
so the family comparison concerns complete controllers rather than an
isolated backbone intervention.

The prespecified screen required at least 0.5\% overall mean cost savings
against H8, positive savings in every training block, a positive descriptive
95\% lower cost bound, and nonpositive mean extra patient losses.
Neither finalist passed (Table~\ref{tab:family-selection}). Value-TD passed
the mean savings and mean patient-loss conditions, but its cost interval
crossed zero and its fifth block cost more. Its five block savings were
0.363, 0.655, 0.401, 0.339, and $-0.072$ million synthetic units.
SAC had higher mean cost in every block and 3.26 additional losses per world.
Thus the frozen rule returned no winner. Value-TD's favorable descriptive
head-to-head contrast with SAC does not override its failed H8 screen.

\begin{table}[tbp]
\centering
\caption{Independent five-block results after development-only family
selection. Savings are reference minus candidate cost, in millions of
synthetic objective units per world; percentages divide the overall mean
saving by the reference mean, not an average of world-level ratios.
Extra losses are candidate minus reference. Intervals are descriptive
95\% intervals, not multiplicity-adjusted inference. The first two rows
are the prespecified H8 screens; the remaining rows separate graph and
training increments. Each row contains 120 paired worlds.}
\label{tab:family-selection}
\small
\setlength{\tabcolsep}{4pt}
\begin{tabularx}{\textwidth}{Xrrr}
\toprule
Candidate / reference & Savings M [95\%] & Savings (\%) & Extra losses \\
\midrule
Value-TD graph / H8 & $0.337\ [-0.026,0.678]$ & $0.851$ & $-11.875$ \\
SAC graph / H8 & $-1.386\ [-2.135,-0.807]$ & $-3.499$ & $3.258$ \\
Value-TD graph / self-only & $0.013\ [-0.115,0.161]$ & $0.033$ & $0.150$ \\
SAC graph / self-only & $-0.261\ [-1.274,0.546]$ & $-0.641$ & $0.525$ \\
Value-TD final / initial & $0.337\ [-0.024,0.676]$ & $0.851$ & $-11.875$ \\
SAC final / initial & $0.005\ [-0.882,0.717]$ & $0.013$ & $3.383$ \\
\bottomrule
\end{tabularx}
\end{table}

All three conditions received equal weight. The 2,000 bootstrap resamples
used common training-block draws across conditions and independent world
draws within each condition; matching replicate labels across conditions
were not treated as paired patients. H8 screen intervals and attribution
intervals retain their frozen contrast-specific random streams. The full
readout retains all 66 development pairs and 36 final pairs, including
condition, block, and world-level harm rather than only aggregate means.
Against H8, value-TD cost more in 51 of 120 worlds and lost more patients
in 21; the corresponding SAC counts were 91 and 75. Value-TD savings
were 0.79\%, 0.25\%, and 1.46\% under stable, persistent, and fast conditions,
but each condition-specific cost interval crossed zero.

Neither graph/self-only comparison established a cost advantage; mean extra
losses were positive for both graph models. Neither final/initial cost
interval excluded zero. The zero-residual value initializer reproduced H8's
cost and loss outcomes, whereas a random SAC initializer is not a strong
operational baseline. These controls do not establish separate message-passing
and training contributions. As secondary references, H16 saved 0.77\%
against H8 with 6.20 fewer mean losses, and the fixed historical controller
saved 0.08\% with 4.79 fewer mean losses; both overall cost intervals crossed
zero. Historical recipe differences remain, and the reference roles were
not newly selected winners.

Mean applied flexible hours per world were 304.13 for H8, 338.73 for H16,
376.86 for value-TD graph, and 328.48 for SAC graph. Mean complete-trajectory
elapsed times were 37.02, 67.98, 36.99, and 4.15 seconds, respectively,
including recorded rollout I/O; these are not separately measured decision
latencies. The complete scientific run, including its automatic analysis and
single archive, took 32.520 hours. Actual consumption was 431,616 native
steps, 445,104 native operations, 1,016,448 neural module forwards,
50,429,952 prediction steps, and 43,161,600 filter transitions.
There were 153,600 minibatch iterations and 316,416 optimizer dispatches:
95,232 actor, 153,600 critic, and 67,584 value steps, presenting
20,250,624 examples. These are actual receipts, not sums of separate budget
ceilings. Saved raw trajectories, summaries, models, and update receipts
support the readout; executing learning is not evidence that learning helped.

The completed experiment therefore establishes neither a qualifying family
winner nor general RL superiority at the declared budget. It does not prove
equivalence or rule out other configurations. No multiplicity-adjusted formal
significance, power, global optimality, real-patient calibration, clinical
safety, deployment-online adaptation, or publication readiness is claimed.
E1 remains missing and all outcomes are synthetic.

\subsection{Post hoc fixed-support MDL-2 comparison}
\label{sec:mdl2-fixed-supplement}
% Evidence: specs/2026-10-08-mdl2-fixed-reference/{protocol,terminal-readout}.md
% Curated numeric binding: reports/2026-10-08-manuscript-publication/evidence.json
The strong-planner comparison leaves a different operational question:
does the learned controller improve on a fixed staffing schedule combined
with the same base operations? After inspecting the family-study outcomes,
we added an MDL-2 fixed-support comparator. Its rule was fixed before its
own outcomes were observed, but the comparator choice was post hoc. This
supplement does not replace the original H8 primary endpoint or its failed
screen, and it must not be confused with the earlier routing-primary MDL-2
comparison.

The fixed rule retains MDL-2 routing and purchasing, ordinary staffing of
four hours per site and epoch, and the original two-epoch delivery lag.
It books two additional support hours at each of four sites during each of
the 48 control epochs, irrespective of the observed queues. No new support
is booked during the 16 settlement epochs. Each world therefore books and
receives 384 flexible hours; idle booked hours, switching, and outstanding
obligations retain their original charges. Feasible resource limits,
physical dynamics, reward, and all three conditions are unchanged. Actual
support consumption is not equalized across policies.

We executed only this new role on the 120 original test tapes
(five blocks, three conditions, eight worlds per condition and block).
The sealed value-TD graph-final, SAC graph-final, and H8 outcomes contribute
360 reused trajectories, with no model loading, neural inference, optimizer
update, or retraining. Full world identities and tape hashes establish
pairing. The new raw records were reconciled with their summaries, while
the previous independently reconciled readout supplies the reused outcomes.
All six pairwise comparisons and their world, block, and condition harms
are retained. The descriptive bootstrap uses 2,000 draws with common
training-block resamples across conditions and independent world resamples
within each condition. Original contrast-specific streams preserve the
previous H8 intervals exactly.

\begin{table}[tbp]
\centering
\caption{Post hoc comparison on 120 reused evaluation worlds. Each entry is
a per-world mean. Cost is in millions of synthetic objective units, not
calibrated expenditure. Savings divide the difference in overall mean costs
by the fixed-support MDL-2 mean. Losses are simulated patient counts;
support is applied additional labor in hours. H8 remains the original
study's primary reference.}
\label{tab:mdl2-fixed-overall}
\small
\setlength{\tabcolsep}{4pt}
\begin{tabularx}{\textwidth}{Xrrrr}
\toprule
Controller & Cost M & Savings (\%) & Losses & Support h \\
\midrule
MDL-2 fixed support & 40.140 & --- & 157.125 & 384.000 \\
H8-MPC & 39.611 & 1.317 & 163.325 & 304.128 \\
Value-TD graph-final & 39.274 & 2.157 & 151.450 & 376.859 \\
SAC graph-final & 40.998 & $-2.136$ & 166.583 & 328.479 \\
\bottomrule
\end{tabularx}
\end{table}

Value-TD saved 0.866 million units per world relative to fixed support
(descriptive 95\% interval $[0.550,1.181]$ million) and lost 5.675 fewer
patients ($[-7.350,-4.000]$ extra losses). All five block-mean cost savings
were positive, but 34 of 120 worlds cost more and 26 lost more patients.
SAC cost 0.857 million more ($[-1.645,-0.278]$ million savings) and lost
9.458 more patients ($[3.075,18.117]$ extra losses); its cost was higher
in every block, with 88 cost-harmed and 104 patient-harmed worlds.
H8 saved 0.529 million ($[0.202,0.892]$ million) but lost 6.200 more
patients ($[3.791,8.642]$ extra losses), with 43 cost-harmed and 70
patient-harmed worlds. Fixed MDL-2 is therefore not uniformly inferior to
H8: its larger labor allocation accompanies a different cost--service
trade-off. The comparison estimates whole-controller differences, not the
causal effect of labor quantity alone.

\begin{table}[tbp]
\centering
\caption{Value-TD graph-final relative to fixed-support MDL-2 by condition,
40 paired worlds per row. Positive savings and negative extra losses favor
value-TD. Intervals are descriptive 95\% bootstrap intervals.}
\label{tab:mdl2-fixed-conditions}
\small
\setlength{\tabcolsep}{3pt}
\begin{tabularx}{\textwidth}{Xrrr}
\toprule
Condition & Savings (\%) & Savings M [95\%] & Extra losses [95\%] \\
\midrule
No change & 3.794 & $1.561\ [0.954,2.167]$ & $-8.875\ [-12.775,-5.350]$ \\
Persistent change & 2.175 & $0.834\ [0.065,1.536]$ & $-6.800\ [-10.275,-3.199]$ \\
Fast fluctuation & 0.494 & $0.202\ [-0.443,0.865]$ & $-1.350\ [-5.150,2.475]$ \\
\bottomrule
\end{tabularx}
\end{table}

The favorable overall means do not establish robustness to fast independent
fluctuations: both its cost and patient-loss intervals span zero
(Table~\ref{tab:mdl2-fixed-conditions}). MDL-2 fixed support required no
planner, response filter, or neural call. Its mean complete-world time was
0.85 seconds, compared with 37.02 for H8, 36.99 for value-TD, and 4.15 for
SAC in their saved runs. These are rollout-plus-I/O measurements, not
synchronized hardware benchmarks or decision-only latency. The supplement
used 7,680 native steps and 7,920 native operations and completed in
112.68 seconds including analysis and archival work. It adds no learning
computation; it does not remove the original models' training cost.

These results support a limited descriptive cost--service advantage of the
tested value-guided controller over fixed allocation in the sampled worlds.
They neither create independent validation from reused tests nor establish
an isolated GCN or RL contribution. The original value-TD/H8 interval still
crosses zero, and the graph/self-only and final/initial cost contrasts remain
inconclusive. Any method revision motivated by these outcomes requires new
development decisions and untouched evaluation worlds, not continued tuning
on this supplement.

\subsection{Prospective fixed-attenuation development comparison}
\label{sec:attenuation-development}
% Evidence: specs/2026-10-08-value-td-attenuation/terminal-readout.md
% Raw reconciliation: reports/2026-10-08-value-td-attenuation/verification.json
We next fixed one development change before observing its outcomes:
multiply only the learned residual in Equation~\ref{eq:current-value-mpc-score}
by 0.5. Prefix cost, terminal heuristic, candidate actions, public predictor,
physical dynamics, and objective weights were unchanged. Three sealed
Value-TD development models at learning rate $10^{-4}$ were reused without
training or checkpoint selection. Each model block used eight fresh worlds
in each of the three conditions. Four roles on these 72 paired worlds
produced 288 trajectories: fixed-support MDL-2, H8, original Value-TD,
and half-residual Value-TD.

The development rule required positive overall savings and no mean loss
increase versus fixed support, together with positive overall savings and
no mean loss increase versus original Value-TD overall or in any condition.
This was an advancement rule, not a formal significance or clinical
noninferiority test. All six contrasts were retained.

\begin{table}[tbp]
\centering
\caption{Development comparison on 72 new worlds using three existing
development models. Savings are millions of synthetic units per world.
Original and Half denote Value-TD with residual multipliers 1 and 0.5.
Intervals are descriptive 95\%; excess losses are candidate minus reference.
All rows have 72 paired observations, conditional on the three models.}
\label{tab:attenuation-pairs}
\small
\setlength{\tabcolsep}{3pt}
\begin{tabularx}{\textwidth}{Xrrr}
\toprule
Candidate / reference & Savings M [95\%] & Savings (\%) & Excess losses \\
\midrule
H8 / fixed MDL-2 & $0.081\ [-0.403,0.548]$ & $0.211$ & $7.111$ \\
Original / fixed MDL-2 & $0.901\ [0.572,1.266]$ & $2.345$ & $-5.181$ \\
Half / fixed MDL-2 & $1.007\ [0.635,1.391]$ & $2.620$ & $-5.681$ \\
Original / H8 & $0.820\ [0.332,1.323]$ & $2.139$ & $-12.292$ \\
Half / H8 & $0.926\ [0.431,1.381]$ & $2.415$ & $-12.792$ \\
Half / Original & $0.106\ [-0.125,0.327]$ & $0.282$ & $-0.500$ \\
\bottomrule
\end{tabularx}
\end{table}

Original Value-TD saved 2.3447\% relative to fixed support and reduced mean
losses by 5.1806 per world (excess-loss interval $[-7.1667,-3.1528]$).
Its cost was worse in 18 of 72 worlds and its loss count was worse in 19.
The half-residual variant's overall means were slightly better, but the
head-to-head cost interval crossed zero and its excess-loss interval was
$[-2.1115,1.4174]$. Against original Value-TD it cost more in 30 worlds
and lost more patients in 28; one model block had increased mean losses.

\begin{table}[tbp]
\centering
\caption{Half-residual minus original behavior by condition. Positive savings
favor attenuation; positive excess losses are unfavorable. Each row contains
24 paired worlds. Harm counts concern cost and patient losses separately.}
\label{tab:attenuation-conditions}
\small
\begin{tabularx}{\textwidth}{Xrrr}
\toprule
Condition & Savings (\%) & Excess losses & Harmed worlds: cost / loss \\
\midrule
No change & $0.4881$ & $-0.7083$ & $7/7$ \\
Persistent unannounced change & $0.7498$ & $0.3750$ & $8/13$ \\
Fast independent fluctuation & $-0.3817$ & $-1.1667$ & $15/8$ \\
\bottomrule
\end{tabularx}
\end{table}

The persistent-change mean loss increase violated the frozen condition-wise
guardrail. The variant therefore did not advance. This is a negative
development decision despite successful execution, not evidence of a
statistically established or clinical harm. Attenuation also did not solve
fast-fluctuation cost performance: versus fixed support, its mean saving
was $-0.1165\%$, with negative cost savings in all three model blocks.
Original Value-TD's corresponding mean saving was 0.2642\%, with cost
interval $[-0.713,0.733]$ million and excess losses 0.1250
($[-2.9167,3.8354]$).

Mean applied additional hours were 384.000 for fixed support, 315.217 for
H8, 377.993 for original Value-TD, and 376.549 for attenuation. Mean full-world
elapsed times including recording were 0.915, 33.879, 33.822, and 33.841
seconds, respectively; these are not isolated decision-latency benchmarks.
The package completed in 7,407.79 seconds, including analysis and archival
work, using 18,432 native steps, 19,008 native operations, 6,912 neural
module forwards, 3,981,312 prediction steps, and 1,382,400 filter transitions.
There were no optimizer updates or clones. Prior model training is not
included in these new execution counts.

This comparison used new development worlds, not new training seeds.
Its favorable original/H8 interval does not supersede the independent
family study's inconclusive original/H8 result. The 72 development worlds
and the earlier 120 test worlds are now consumed for future confirmation.
No second coefficient was selected or evaluated.

\section{Discussion and Conclusions}
\label{sec:conclusion}
The main finding is comparator-dependent. Original Value-TD shows favorable
mean cost and simulated patient losses relative to the specified fixed
support schedule, first in a post hoc supplement and then on new
development worlds conditional on three existing models. These findings
do not establish a qualifying winner against strong MPC under the original
five-block screen. H8 uses less support and may trade lower cost against
higher patient losses. It remains a necessary reference rather than an
unfavorable result to omit.

The rejected attenuation study also separates mean improvement from
development acceptance. Scaling the learned residual by one half modestly
improved aggregate means, but failed the predeclared condition-wise patient
rule and did not resolve fast-fluctuation cost weakness. The original
controller is retained as a research candidate. The unsuccessful variant
is reported rather than used to justify further coefficient searches.

The method's distinctive content lies in the support-hour formulation and
the explicit integration of fixed operational control, approximate public
beliefs, graph-valued remaining cost, and settlement-complete evaluation.
This is not yet evidence of a new general-purpose algorithm. A new
application alone does not guarantee sufficient methodological novelty for
publication. Graph/self-only and final/initial controls have not established
separate cost benefits; the title's graph-aware description names the
implementation, not a demonstrated causal graph gain.

Three limits are central. First, public-state compression and forecast error
can affect candidate ranking. Second, simulated work response, staffing
notice, costs, and eligibility are not externally calibrated. Published
process evidence motivates assumptions but cannot validate numerical
operating savings or clinical benefits. Third, favorable development
contrasts after method selection do not supply independent confirmation
across fresh training seeds or unseen networks. More resampling of the same
blocks cannot create those independent units.

The next empirical question is therefore bounded: with the unchanged
original recipe fixed in advance, do fresh training and untouched evaluation
support improvement over MDL-2 fixed support, and what increments remain
against H8 and a matched self-only value model? A corresponding protocol
must prospectively specify seeds, patient guardrails, uncertainty, complete
compute charges, and interpretation when costs and losses trade off.
The zero-residual initializer shares H8's planner scores; a matched initial
control should be accounted for explicitly rather than mislabeled as an
additional strong baseline. No such follow-up has been executed here.

Separately, engineering validity requires a process-specific mapping from
epochs and preparatory work to observed tasks, qualified staff availability,
booking lead times, and an economic boundary. Until that evidence exists,
a small justified sensitivity panel can assess synthetic robustness, not
replace calibration. An intervention that changes physical transport or
sharing must not be called a neural graph ablation.

We conclude that the current evidence warrants continued evaluation of
Value-TD as a synthetic decision-support candidate, not deployment,
universal dominance, isolated graph/RL attribution, or a publication
guarantee. Independent confirmation and an accountable assessment of the
engineering assumptions are the next research requirements.

\section*{Data and Code Availability}
All outcomes in this manuscript were generated by the simulation model.
The local research record contains locked configurations, raw trajectories,
model and update provenance, paired summaries, and SHA256 archive manifests.
The three support-hour readouts are identified in the accompanying
claim-to-evidence map. Historical full-routing and earlier value-learning
studies are preserved separately and are not pooled with the current results.
Code and earlier curated records have been published at
\texttt{github.com/ZhaoweiLi0722/GCN-RL-Paper-2026}; this revision and its latest
development artifacts are local, with no new remote release or collaborator
access verification claimed. Final availability wording requires author review.

\section*{Writing Assistance and Declarations}
This working revision used Codex and the scientific-writing and
manuscript-writing-review instructions for drafting and consistency checks;
the Scientific Agent Skills reference records procedural writing
assistance~\cite{Kassis2026ScientificAgentSkills}, not evidence for the
controller. Human authors must verify sources, scientific claims,
authorship, funding and competing-interest declarations, and the final
disclosure and submission decision. No author approval is inferred here.

\bibliography{references}
\end{document}
